% Evans 12.2, printed pp.620--627; English PDF pp.637--644.
% Chinese start: printed p.439 / scan p.457.
\Needspace{45mm}
\section{Existence of Solutions}
\label{sec:12-2}
\subsection{The linear estimate and the energy solution class}
For an integer $k\ge1$, write
\[
 \|(u,u_t)\|_{\mathcal E_k(I)}
 =\sup_{t\in I}\{\|u(t)\|_{H^k}+\|u_t(t)\|_{H^{k-1}}\}.
\]
The inhomogeneous $H^k$ norm includes the low-frequency $L^2$ term.
We use the whole-space linear wave construction from Chapters 2 and 7:
Fourier multipliers $\cos(t|D|)$ and $\sin(t|D|)/|D|$ give, for
$u_{tt}-\Delta u=G$,
\begin{equation}\label{eq:12-linear-energy}
 \|(u,u_t)\|_{\mathcal E_k([0,T])}
 \le C_T\left(\|g\|_{H^k}+\|h\|_{H^{k-1}}
                    +\|G\|_{L^1(0,T;H^{k-1})}\right).
\end{equation}
Here $C_T$ is bounded for $T\le1$.
\begin{expandedproof}[the energy estimate with its missing low mode]
For smooth data apply $D^\alpha$, $|\alpha|\le k-1$, and multiply
by $D^\alpha u_t$. The square root of
$\|D^\alpha u_t\|_2^2+\|DD^\alpha u\|_2^2$ grows at most by
$\|D^\alpha G\|_2$. Replace this square root by
$\sqrt{\varepsilon+\cdots}$ if it vanishes, then let $\varepsilon\downarrow0$.
The remaining $\|u(t)\|_2$ is bounded by
$\|g\|_2+\int_0^t\|u_t(s)\|_2ds$.
Summing gives \eqref{eq:12-linear-energy}; density gives a solution
in $C_tH^k\cap C_t^1H^{k-1}$. Duhamel's formula gives continuity
also for $G\in L^1_tH^{k-1}$. Thus no Poincare constant depending
on an expanding spatial ball is needed.
\end{expandedproof}

An energy solution below has $u\in C([0,T];H^1)$,
$u_t\in C([0,T];L^2)$ and satisfies the PDE in distributions.
If $N(Du,u_t,u)\in L^2_{\rm loc}$, then
$u_{tt}\in L^2_{\rm loc}(0,T;H^{-1}_{\rm loc})$ and the usual
weak formulation against compactly supported $H^1$ tests follows.
The two time traces specify $g,h$ separately.

\subsection{Globally Lipschitz nonlinearities}
\begin{theorem}[Source \S12.2.1, Theorem 1]
Assume $N:\R^{n+2}\to\R$ is globally Lipschitz, $N(0)=0$.
For $g\in H^1$, $h\in L^2$ there is a unique global energy solution
of $\PDEWave u+N(Du,u_t,u)=0$. For $g\in H^2$, $h\in H^1$ it also has
$u\in L^\infty(0,T;H^2)$, $u_t\in L^\infty(0,T;H^1)$,
$u_{tt}\in L^\infty(0,T;L^2)$ on every finite strip.
With $H^1_{\rm loc},L^2_{\rm loc}$ data the corresponding statements
hold locally in space, with uniqueness in the local energy class.
\end{theorem}
\begin{expandedproof}[contraction, restart, and local gluing]
On the complete space $C_tH^1\cap C_t^1L^2$ with the prescribed
initial traces, define $A[v]$ by the linear equation with forcing
$-N(Dv,v_t,v)$. Since $N(0)=0$, this forcing belongs to $C_tL^2$.
For two inputs, \eqref{eq:12-linear-energy} gives
\[
 \|A[v]-A[w]\|_{\mathcal E_1}
 \le C_1LT\|v-w\|_{\mathcal E_1}.
\]
The homogeneous solution makes the affine space nonempty. Choose
$T\le1$ so $C_1LT<1/2$ and apply Banach's theorem.
The same step size works for every new initial state, so repeated
restart constructs a solution on every finite interval. Its norm is
finite there, for example by the integral estimate and Gronwall;
no finite accumulation of restart times occurs.

For the additional derivative, apply the spatial difference quotient
$\delta_\varepsilon$ to the PDE. Pointwise Lipschitz continuity gives
\[
 |\delta_\varepsilon N(Du,u_t,u)|
 \le L(|D\delta_\varepsilon u|+|\delta_\varepsilon u_t|
                                  +|\delta_\varepsilon u|).
\]
The difference energy and Gronwall bound the quotient in
$L^\infty_tH^1$ and its time derivative in $L^\infty_tL^2$,
uniformly in $\varepsilon$, by the $H^2\times H^1$ initial norm.
Weak compactness and the difference-quotient characterization give
$D^2u,Du_t\in L^\infty_tL^2$; the equation then gives the asserted
bound for $u_{tt}$. Differentiability of a merely Lipschitz $N$ is
not needed for this argument.

For local data, cut them off outside $B(0,j)$ and solve the global
problem. Two cutoff solutions coincide in the backward domains where
their data agree, by the regularized difference-cone estimate from
\S\ref{sec:12-1}. For every fixed compact space-time cylinder they
therefore agree for all sufficiently large $j$. Their common values
define a solution independent of the cutoff. The same cone argument
proves uniqueness and localizes the higher-regularity estimates.
\end{expandedproof}
\PDEAnnotation{原点条件与全空间}{若 $N(0)\ne0$，零状态也会受到常量
强迫，$N(Dv,v_t,v)$ 不一定在 $L^2(\R^n)$。
$N(0)=0$ 同时保证零解、有限传播中的比较基准与全空间能量构造。}

\subsection{Sobolev products and composition}
\begin{theorem}[Source \S12.2.2, Theorem 2; Problem 12]
Let $r>n/2$ be an integer, $u_1,\ldots,u_m\in H^r(\R^n)$.
For multiindices with $\sum_j|\beta_j|\le r$,
\begin{equation}\label{eq:12-product}
 \left\|\prod_jD^{\beta_j}u_j\right\|_2
 \le C_{n,r,m}\prod_j\|u_j\|_{H^r}.
\end{equation}
If $P:\R^m\to\R$ is smooth and $P(0)=0$, then $P(u)\in H^r$,
with a bound depending on the $H^r$ norms and derivatives of $P$
on the resulting bounded range.
\end{theorem}
\begin{expandedproof}[including the exponent allocation]
Put $b_j=|\beta_j|$. Sobolev embedding gives
$D^{\beta_j}u_j\in L^{2r/b_j}$ for $0<b_j\le r$ and $u_j\in L^\infty$
for $b_j=0$. Indeed for $b_j<r$ the choice $1/p=b_j/(2r)$ lies
strictly above $1/2-(r-b_j)/n$, since $2r>n$; $b_j=r$ gives $p=2$.
Each derivative also lies in $L^2$. By interpolation choose
$1/p_j\in[b_j/(2r),1/2]$ so that $\sum_j1/p_j=1/2$.
This is possible because the sum of the lower endpoints is at most
$1/2$ and that of the upper endpoints is at least $1/2$.
Holder's inequality proves \eqref{eq:12-product}. On the infinite
measure space one must make this sum exact, rather than infer an
$L^2$ estimate from a product lying only in a larger $L^p$ space.

For $0<|\alpha|\le r$, the chain rule expresses $D^\alpha P(u)$
as a finite sum of bounded coefficients $D^jP(u)$ times products
of derivatives of the $u_i$ with total order $|\alpha|$.
Apply \eqref{eq:12-product}. For the undifferentiated term use
$P(0)=0$ and $|P(u)|\le C|u|$ on the bounded range.
Approximation by smooth functions proves the Sobolev chain rule
and the assertion for all $H^r$ inputs. The same estimates applied
to $P(u)-P(v)=\int_0^1DP(v+\theta(u-v))(u-v)d\theta$ show local
Lipschitz dependence in $H^r$; subtract the constant $DP(0)$ before
using the algebra estimate on its coefficient.
\end{expandedproof}

\subsection{Local existence and the role of the two norms}
\begin{theorem}[Source \S12.2.2, Theorem 3]
Assume $N$ is smooth with $N(0)=0$, and $k>n/2+1$ is an integer.
For $g\in H^k$, $h\in H^{k-1}$ there are $T>0$ and a unique solution
\[
 u\in C([0,T];H^k),\qquad u_t\in C([0,T];H^{k-1}).
\]
The lifetime depends on these norms and the nonlinearity on a bounded
state range. Classical $C^2$ regularity follows if $k>n/2+2$.
\end{theorem}
\PDEAnnotation{两个范数各做什么}{高阶界使 $(Du,u_t,u)$ 的值域有界，
从而给出 $N$ 的局部 Lipschitz 常数。低阶能量范数用于收缩。
只说迭代在弱范数收敛，还需证明强范数球在该度量下闭合。}
\begin{expandedproof}[closed strong ball with the weaker metric]
Fix $T\le1$ and let $Y_R$ consist of pairs $(v,v_t)$ continuous in
$H^1\times L^2$, with the prescribed traces, and with
$\|(v,v_t)\|_{\mathcal E_k}\le R$ in the essential-supremum sense.
Use the $\mathcal E_1$ metric. This space is complete: a Cauchy
sequence has a strong low-norm limit; the high-norm bounded sequence
has a weak-star subsequence in $L^\infty(H^k\times H^{k-1})$.
Both limits agree as distributions. Weak lower semicontinuity gives
the bound $R$, and low-norm convergence preserves the derivative
relation and both initial traces. The uniform high bound can also
be assigned at each time by weak limits of nearby good time slices.

For $A[v]$ as above, composition at order $k-1>n/2$ gives
\[
 \|N(Dv,v_t,v)\|_{H^{k-1}}\le\Phi(R).
\]
Take $R>2C_1(\|g\|_{H^k}+\|h\|_{H^{k-1}})+1$.
The free wave belongs to $Y_R$, and
\[
 \|A[v]\|_{\mathcal E_k}\le C_1D_0+C_1T\Phi(R)\le R
\]
for small $T$. Sobolev embedding bounds all first-order states by
$C R$, so the local Lipschitz constant $L_R$ gives
$\|A[v]-A[w]\|_{\mathcal E_1}\le C_1TL_R\|v-w\|_{\mathcal E_1}$.
Shorten $T$ to make this less than one half of the input distance.
Banach's theorem gives a fixed point in the closed space $Y_R$.
The resulting forcing is in $L^1_tH^{k-1}$; linear uniqueness and
\eqref{eq:12-linear-energy} then upgrade the solution to the asserted
strong time continuity. Differences of two bounded high-regularity
solutions obey the same low-energy Gronwall estimate, proving uniqueness
without restricting them to one preselected ball. If $k>n/2+2$,
$D^2u,Du_t,u_{tt}$ are continuous by embedding and the PDE.
\end{expandedproof}
\PDEBodyDiagram{12.2}{A complete ball for the weak-metric iteration}
{Every iterate stays in a bounded high Sobolev ball. Contraction is measured
in the lower energy norm; weak-star compactness identifies its limit with
the low-norm limit, preserving the high bound. The final linear equation
restores strong time continuity. The arrows describe a proof, not a
numerical convergence rate in the high norm.}{ch12-two-norms}
